% AI: Combine: less text than in version 1 but more formulas than in version 2

\section{Gradient Methods - Part II}

\subsection{Projection and Projected Gradient Descent}
For a closed convex set $\mathcal{C}$ the projection is
\[
  \operatorname{prox}_{\mathcal{C}}(x)=\arg\min_{y\in\mathcal{C}}\frac12\|x-y\|^2.
\]
It is firmly non-expansive:
\begin{lemma}
  $|\operatorname{prox}_{\mathcal{C}}(x)-\operatorname{prox}_{\mathcal{C}}(y)|^2\le(\operatorname{prox}_{\mathcal{C}}(x)-\operatorname{prox}_{\mathcal{C}}(y))^\top(x-y)$,
  implying
  $\|\operatorname{prox}_{\mathcal{C}}(x)-\operatorname{prox}_{\mathcal{C}}(y)\|\le\|x-y\|$.
\end{lemma}

The projected gradient step is
$x_{k+1}=\operatorname{prox}_{\mathcal{C}}(x_k-T\nabla f(x_k))$. For
$L$-smooth $\mu$-strongly convex $f$ and $T=2/(L+\mu)$ the same
linear rate as the unconstrained case holds:
\[
  \|x_N-x^\star\|\le\Bigl(1-\frac{2}{\kappa+1}\Bigr)^N\|x_0-x^\star\|.
\]
Proof (Lecture 6): the fixed point property
$x^\star=\operatorname{prox}_{\mathcal{C}}(x^\star-T\nabla
f(x^\star))$ follows from the optimality condition $T\nabla
f(x^\star)+\partial\psi_{\mathcal{C}}(x^\star)\ni0$; the contraction
follows by non-expansiveness of the prox and the same contraction of
$I-TH$ as in the unconstrained quadratic case.

Projections are cheap for Euclidean balls
($\operatorname{prox}(y)=y/\max(1,\|y\|)$), $\ell_1$-balls
(soft-thresholding after sorting), probability simplices, and
low-dimensional affine subspaces.

\subsection{Smooth but Not Strongly Convex ($\mu=0$)}
Add a small strongly convex regularizer:
\[
  \tilde{f}(x)=f(x)+\frac{\mu}{2}\|x-x_0\|^2.
\]
Lemma (Lecture 6): $\tilde{f}$ is $(L+\mu)$-smooth and $\mu$-strongly
convex,
$f(x)-f(x^\star)\le\tilde{f}(x)-\tilde{f}(\tilde{x}^\star)+(\mu/2)\|x^\star-x_0\|^2$,
and $\|\tilde{x}^\star-x_0\|\le\|x^\star-x_0\|$.

Applying Nesterov's accelerated method to $\tilde{f}$ and optimizing
the regularization parameter $\mu\approx(2L\ln N)/N$ yields (after
using $e^\xi\ge1+\xi$)
\[
  f(x_N)-f(x^\star)\le\Bigl(\frac{2L}{N^2}+\frac{9L\ln(N)^2}{2N^2}\Bigr)\|x^\star-x_0\|^2.
\]
A refined analysis gives the optimal accelerated rate (Beck–Teboulle FISTA):
\begin{theorem}
  The scheme $q_{k+1}=q_k+p_{k+1}$, $p_{k+1}=\beta_k
  p_k-g_f(q_k+\beta_k p_k)/L$ with $\beta_k=t_k/t_{k+1}$,
  $t_{k+1}=(1+\sqrt{1+4t_k^2})/2$, $t_0=1$ satisfies
  \[
    f(q_N)-f(x^\star)\le\frac{2L\|q_0-x^\star\|^2}{(N+1)^2}.
  \]
\end{theorem}
A matching lower bound of order $1/N^2$ exists.

Thus, roughly $N\sim\sqrt{L\|x^\star-x_0\|^2/\varepsilon}$ iterations
suffice for $\varepsilon$-accuracy.

\subsection{Non-Smooth Convex Objectives}
If $f$ is convex and $L_f$-Lipschitz (but not differentiable),
constant-step gradient descent can oscillate (example $f(x)=|x|$).
Use a diminishing stepsize. The projected subgradient method with
$T=R L_f/\sqrt{N}$ on a set of radius $R$ yields
\[
  f\Bigl(\frac1N\sum_{k=0}^{N-1}x_k\Bigr)-f(x^\star)\le\frac{R L_f}{\sqrt{N}}.
\]
Proof (Lecture 6): combine the subgradient inequality $f(x^\star)\ge
f(x_k)+g_k^\top(x^\star-x_k)$, the non-expansiveness of the prox
(distance to $x^\star$ decreases), and telescope; convexity of $f$
lets us bound the average.

\begin{table}[ht]
  \centering
  \begin{tabular}{llll}
    \toprule
    Assumptions on $f$ & Method & $N$ for
    $f(x_N)-f(x^\star)\le\varepsilon$ & Optimal? \\
    \midrule
    $\mu$-strongly convex, $L$-smooth & GD / Nesterov &
    $\sim\kappa\ln(1/\varepsilon)$ /
    $\sim\sqrt{\kappa}\ln(1/\varepsilon)$ & Yes (Nesterov) \\
    $L$-smooth convex & Nesterov (accelerated) &
    $\sim1/\sqrt{\varepsilon}$ & Yes \\
    $L_f$-Lipschitz, compact $\mathcal{C}$ & Subgradient &
    $\sim1/\varepsilon^2$ & Yes \\
    \bottomrule
  \end{tabular}
\end{table}

The plot of error vs.\ iteration (Lecture 6) clearly shows linear,
$1/k$, $1/k^2$, and $1/\sqrt{k}$ regimes.

Recitation remarks on adaptive stepsize rules, normalized gradient
descent, and practical tuning close the algorithmic part. These
rates, together with the theoretical foundations of the first four
chapters, provide a complete toolkit for recognizing, formulating,
and solving large-scale convex problems with first-order methods.

\section{Gradient Methods – Part II}

\subsection{Projected Gradient Descent}
Let
\(\operatorname{prox}_{\mathcal{C}}(x)=\arg\min_{y\in\mathcal{C}}\frac12\|x-y\|^2\).
For a closed convex set the proximal operator is firmly non-expansive:
\[
  \|\operatorname{prox}_{\mathcal{C}}(x)-\operatorname{prox}_{\mathcal{C}}(y)\|\le\|x-y\|.
\]
The iteration \(x_{k+1}=\operatorname{prox}_{\mathcal{C}}(x_k-T\nabla
f(x_k))\) with \(T=2/(L+\mu)\) converges exactly as the unconstrained
method when \(f\) is smooth and strongly convex.

\textbf{Cheap projections.} Euclidean ball, \(\ell_1\)-ball,
probability simplex, affine subspaces.

\subsection{Non-Strongly-Convex Case}
Regularise:
\[
  \tilde f(x)=f(x)+\tfrac\mu2\|x-x_0\|^2.
\]
Apply accelerated gradient descent to \(\tilde f\) and optimise the
regularisation parameter \(\mu\sim L\log N/N\). The resulting sub-linear rate is
\[
  f(x_N)-f(x^*)\le O\Bigl(\frac{L\|x_0-x^*\|^2}{N}\Bigr).
\]
Hence \(N\sim L\|x_0-x^*\|^2/\varepsilon\) iterations are necessary
and sufficient.

\subsection{Non-Smooth Case}
When \(f\) is merely Lipschitz (subgradient method) the best
achievable rate on a bounded domain is \(O(1/\sqrt N)\). This rate is optimal.

\clearpage

